\section{Conclusion}
\label{sec:conclusion}

FraudShield-EAC was built to develop and evaluate a fraud detection system for East African digital
payments without real customer data. Measured against it, the specified ensemble ranks fraud well,
but the benchmark sets a high floor: one velocity feature alone separates the classes at
\nFloorM{}, and the ensemble's margin over the floor on the same rows is \nGateMargin{}. Most of what
a reader would like to learn from the benchmark it cannot teach. Feature ablations measure redundancy,
because four disjoint groups each nearly solve the task; leave-one-country-out and a novel variant
held out in time are null, because every country and variant is a view of one burst structure; and
the one variant pre-registered without that structure is caught at \nRevRecall{}. What survives is a
methodology, the record of how each of these limits was found, and a cost that does not depend on how
easy the data is: the price of an exact explanation grows with the model while accuracy does not.

\paragraph{Future work.} Four pieces of work would turn these findings into claims about real payments.
First, validation on real, labelled transactions, under a data protection impact assessment and in a
permitted location, measuring the same floor, margins and calibration. Second, a held-out variant with
a new recipient, as in the documented form of the reversal scam, which would separate the contributions
of burst structure and counterparty novelty that this paper's variant removes together. Third, the
learning curve over training volume, which will say whether the recall shortfall is a property of the
benchmark or of the \nTrainUsed{} training rows: \notyet{PB-67}. Fourth, the serving and load
measurements on dedicated hardware: the frontier at the 99th percentile, end-to-end decision latency,
and scale to five million rows: \notyet{M5, M6 and M10}.
